% Options for packages loaded elsewhere
\PassOptionsToPackage{unicode}{hyperref}
\PassOptionsToPackage{hyphens}{url}
\documentclass[
  a4paper,
]{article}
\usepackage{xcolor}
\usepackage[margin=25mm]{geometry}
\usepackage{amsmath,amssymb}
\setcounter{secnumdepth}{-\maxdimen} % remove section numbering
\usepackage{iftex}
\ifPDFTeX
  \usepackage[T1]{fontenc}
  \usepackage[utf8]{inputenc}
  \usepackage{textcomp} % provide euro and other symbols
\else % if luatex or xetex
  \usepackage{unicode-math} % this also loads fontspec
  \defaultfontfeatures{Scale=MatchLowercase}
  \defaultfontfeatures[\rmfamily]{Ligatures=TeX,Scale=1}
\fi
\usepackage{lmodern}
\ifPDFTeX\else
  % xetex/luatex font selection
\fi
% Use upquote if available, for straight quotes in verbatim environments
\IfFileExists{upquote.sty}{\usepackage{upquote}}{}
\IfFileExists{microtype.sty}{% use microtype if available
  \usepackage[]{microtype}
  \UseMicrotypeSet[protrusion]{basicmath} % disable protrusion for tt fonts
}{}
\makeatletter
\@ifundefined{KOMAClassName}{% if non-KOMA class
  \IfFileExists{parskip.sty}{%
    \usepackage{parskip}
  }{% else
    \setlength{\parindent}{0pt}
    \setlength{\parskip}{6pt plus 2pt minus 1pt}}
}{% if KOMA class
  \KOMAoptions{parskip=half}}
\makeatother
\usepackage{longtable,booktabs,array}
\newcounter{none} % for unnumbered tables
\usepackage{calc} % for calculating minipage widths
% Correct order of tables after \paragraph or \subparagraph
\usepackage{etoolbox}
\makeatletter
\patchcmd\longtable{\par}{\if@noskipsec\mbox{}\fi\par}{}{}
\makeatother
% Allow footnotes in longtable head/foot
\IfFileExists{footnotehyper.sty}{\usepackage{footnotehyper}}{\usepackage{footnote}}
\makesavenoteenv{longtable}
\usepackage{graphicx}
\makeatletter
\newsavebox\pandoc@box
\newcommand*\pandocbounded[1]{% scales image to fit in text height/width
  \sbox\pandoc@box{#1}%
  \Gscale@div\@tempa{\textheight}{\dimexpr\ht\pandoc@box+\dp\pandoc@box\relax}%
  \Gscale@div\@tempb{\linewidth}{\wd\pandoc@box}%
  \ifdim\@tempb\p@<\@tempa\p@\let\@tempa\@tempb\fi% select the smaller of both
  \ifdim\@tempa\p@<\p@\scalebox{\@tempa}{\usebox\pandoc@box}%
  \else\usebox{\pandoc@box}%
  \fi%
}
% Set default figure placement to htbp
\def\fps@figure{htbp}
\makeatother
\setlength{\emergencystretch}{3em} % prevent overfull lines
\providecommand{\tightlist}{%
  \setlength{\itemsep}{0pt}\setlength{\parskip}{0pt}}
\usepackage{bookmark}
\IfFileExists{xurl.sty}{\usepackage{xurl}}{} % add URL line breaks if available
\urlstyle{same}
\hypersetup{
  hidelinks,
  pdfcreator={LaTeX via pandoc}}

\author{}
\date{}

\begin{document}

\section{Paper 6}\label{paper-6}

\section{Arrest of Spontaneous Splitting and Emergence of New Orthogonal
Rotation
Planes}\label{arrest-of-spontaneous-splitting-and-emergence-of-new-orthogonal-rotation-planes}

\textbf{Noriaki Kihara}\\
\textbf{2026-07-24}\\
\textbf{DOI:} 10.5281/zenodo.21543071 (Concept DOI:
10.5281/zenodo.21543070)

\subsection{Abstract}\label{abstract}

In earlier papers we started from the two simple axioms

\[
\sum_n x_n^2=0,
\qquad
U^n=I
\]

and showed that, even in a closed system that assumes no background
spacetime, three directions readable as spatial directions and one
direction readable as a conserved quantity arise spontaneously
{[}1,2{]}.

We further showed that a wave initially present as a single mode is
amplified geometrically from a seed too small to measure, and
spontaneously splits into several waves. This amplification is an
inflation-like exponential expansion, but it does not continue
indefinitely; it moves into a metastable state of finite amplitude
{[}3{]}.

It remained open, however, why the spontaneous splitting stops, and what
happens inside the system at the transition to the metastable state.

In this paper we measure the arrest process directly. We find that, at
the moment spontaneous splitting stops and the system moves into the
metastable state, several orthogonal rotational directions that did not
exist before, together with a new direction readable as a conserved
quantity, arise spontaneously.

The claim of this paper is therefore:

\[
\boxed{
\text{Spontaneous splitting of a single wave stops together with the creation of new orthogonal directions.}
}
\]

The waves produced by splitting are not accumulated into the single
existing rotational direction. At the arrest point of splitting, new
mutually orthogonal rotational directions stand up, and the conserved
quantity is distributed into those directions. The metastable state is
not a state where splitting has merely weakened; it is the state after
the system has acquired new directions.

With this result, the change observed in the earlier paper,

\[
\text{exponential spontaneous splitting}
\longrightarrow
\text{metastable state},
\]

can be described as

\[
\boxed{
\text{localization in a single direction}
\longrightarrow
\text{spontaneous splitting}
\longrightarrow
\text{creation of new orthogonal directions}
\longrightarrow
\text{metastable state with multiple directions}
}.
\]

In the metastable state, the amplitude per relational wave follows

\[
A_{\mathrm{rel}}
\simeq
\frac{1}{\sqrt{M}}
=
\sqrt{\frac{2}{N(N-1)}}.
\]

This means that as the system size \(N\) grows, the amplitude of each
relational wave decreases in proportion to \(1/N\). The reason the
metastable amplitude appears dual to \(N\) is that the conserved
quantity is distributed over

\[
M=\frac{N(N-1)}{2}
\]

relational waves.

The creation of new orthogonal rotational directions, its simultaneity
with the arrest of splitting, the inflow of the conserved quantity, and
the \(1/\sqrt{M}\) law of the amplitude presented here are results
obtained from the dynamics and numerical experiments of the model; they
are not hypotheses.

What is currently undetermined is what the newly arising directions
correspond to in actual physics. They might be identifiable as a time
direction, charge directions, or other internal degrees of freedom, but
that physical identification is not part of the claim of this paper.
Here we show only the existence and creation mechanism of the new
directions, and leave their naming and physical meaning to future work.

\begin{center}\rule{0.5\linewidth}{0.5pt}\end{center}

\subsection{1. System and observables}\label{system-and-observables}

\subsubsection{1.1 State space}\label{state-space}

We represent the complete pairwise relations of \(N\) vertices as edges,
with edge number

\[
M=\binom{N}{2}=\frac{N(N-1)}{2},
\]

and the state as

\[
Z\in\mathbb C^M.
\]

We use the two quadratic quantities

\[
Z^\dagger Z=1,
\qquad
Q(Z)=Z^T Z=0.
\]

Writing \(Z=X+iY\), the zero-square closure is equivalent to

\[
\|X\|=\|Y\|,
\qquad
X\cdot Y=0,
\]

so the state is represented by orthogonal, equal-length components in
real \(M\)-dimensional space.

\subsubsection{1.2 Generator}\label{generator}

From the edge phases

\[
\theta=\arg Z,
\]

we construct, via the low-rank vertex decomposition derived earlier
{[}1{]}, the real antisymmetric generator

\[
K(\theta)=WJW^T,
\qquad
K^T=-K.
\]

Time evolution is the Cayley transform of the normalized generator
\(\widetilde K=K/\sigma_1\),

\[
Z_{\tau+1}
=
\left(I-\gamma\widetilde K_\tau\right)^{-1}
\left(I+\gamma\widetilde K_\tau\right)Z_\tau,
\qquad
\gamma=\tan\frac{\pi}{144}.
\]

Each Cayley matrix is a real orthogonal matrix, so it preserves
\(Z^\dagger Z\) and \(Z^T Z\). In the numerical runs the deviation of
the zero-square closure was at most about \(10^{-14}\).

\subsubsection{1.3 Parent state and dormant
seed}\label{parent-state-and-dormant-seed}

The initial parent state \(v\) is constructed, following the plane
decomposition of the fixed-generator system {[}2{]}, as the dominant
eigenmode of the generator built from its own phases,

\[
K(\arg v)v=-i\sigma_1 v,
\]

so the parent is localized in a single real two-dimensional rotation
plane

\[
P_1=\operatorname{span}\{\Re v,\Im v\}.
\]

To this parent we add a small zero-closure seed \(g\) lying in the
kernel of the initial generator:

\[
Z_0=\frac{v+\delta g}{\|v+\delta g\|},
\qquad
K(\arg v)g=0,
\qquad
\delta=10^{-15}.
\]

\subsubsection{1.4 Splitting quantity}\label{splitting-quantity}

We define the projection ratio onto the parent plane,

\[
P_1(\tau)
=
\frac{|p^T Z_\tau|^2+|q^T Z_\tau|^2}{Z_\tau^\dagger Z_\tau},
\]

and the fraction outside the parent plane,

\[
f(\tau)=1-P_1(\tau),
\]

where \(p,q\) are an orthonormal basis of the parent plane.

\begin{center}\rule{0.5\linewidth}{0.5pt}\end{center}

\subsection{2. Onset mechanism of
splitting}\label{onset-mechanism-of-splitting}

\subsubsection{2.1 State-dependent
generator}\label{state-dependent-generator}

Under the fixed generator \(K_0=K(\arg v)\), the kernel component \(g\)
satisfies \(K_0g=0\) and does not grow.

In the original dynamics, however, the generator
\(K_\tau=K(\arg Z_\tau)\) is reconstructed at each step. When the small
seed changes the state phases, the generator also changes, and a
component that was in the kernel initially ceases to be in the kernel of
the next generator. This feedback amplifies the out-of-parent-plane
component.

\subsubsection{2.2 Generator-freezing
experiment}\label{generator-freezing-experiment}

For \(N=6\) we compared two runs with the same parent state and same
seed:

\begin{itemize}
\tightlist
\item
  fixed generator: \(K_\tau=K_0\)
\item
  updated generator: \(K_\tau=K(\arg Z_\tau)\)
\end{itemize}

With the fixed generator, \(f\) stayed at about \(10^{-27}\) and did not
grow for 2000 steps. With the updated generator, an amplification of
about 30 orders of magnitude arose from the same initial seed, reaching
\(f\simeq0.70\).

Thus the amplification source of splitting is not the action of a fixed
generator, but the feedback between the state phases and the generator.

In the initial amplification regime \(f(\tau)\simeq f_0 e^{r\tau}\)
holds; for \(N=6\) we obtained \(r\simeq0.044\), consistent with
\(r\simeq2\lambda_\perp\) where \(\lambda_\perp\) is the transverse
perturbation growth rate.

\subsubsection{2.3 Conclusion}\label{conclusion}

We call this mechanism \textbf{self-excited amplification by phase
feedback}. Since the total norm is conserved, the amplification is not
the creation of norm but a redistribution from the parent plane to its
orthogonal complement.

\begin{center}\rule{0.5\linewidth}{0.5pt}\end{center}

\subsection{3. Arrest of splitting and
metastabilization}\label{arrest-of-splitting-and-metastabilization}

\subsubsection{3.1 Not a simple depletion of parent
norm}\label{not-a-simple-depletion-of-parent-norm}

Let \(f_{\mathrm{sat}}\) denote the metastable value of the splitting
quantity. Representative values:

{\def\LTcaptype{none} % do not increment counter
\begin{longtable}[]{@{}rrr@{}}
\toprule\noalign{}
N & f\_sat & parent-plane residual 1-f\_sat \\
\midrule\noalign{}
\endhead
\bottomrule\noalign{}
\endlastfoot
5 & 0.80 & 0.20 \\
40 & 0.20 & 0.80 \\
300 & 0.083 & 0.917 \\
\end{longtable}
}

For \(N=40\) and \(N=300\), the splitting quantity saturates while most
of the norm remains in the parent plane. Therefore the arrest cannot be
explained by depletion of the parent-plane norm alone.

What is established at present is that

\[
\boxed{f_{\mathrm{sat}}(N)\text{ decreases as }N\text{ increases}}
\]

and that this arrest is not a simple resource depletion. The closed form
and geometric origin of \(f_{\mathrm{sat}}(N)\) are not yet derived.

\subsubsection{3.2 The metastable state is not a fixed
point}\label{the-metastable-state-is-not-a-fixed-point}

After saturation the state continues to change in time. For
\(N=5\)--\(7\), the transverse growth rate in the metastable regime
remained about 80\% of its initial value. Hence metastabilization is not
convergence to a stable fixed point.

By the metastable state here we mean a bounded non-stationary state with

\begin{itemize}
\tightlist
\item
  \(Z^\dagger Z=1\)
\item
  \(Z^T Z=0\)
\item
  \(f\) staying in a finite interval
\item
  the state itself continuing to fluctuate.
\end{itemize}

A rigorous test of ``mixing'' would require additional measurements such
as the Lyapunov spectrum, autocorrelation, and recurrence statistics.

\subsubsection{3.3 Spectral indicator}\label{spectral-indicator}

Let the nonzero rotation rates of the generator be
\(\sigma_1\ge\sigma_2\ge\cdots>0\). The second-rate ratio
\(\rho_2=\sigma_2/\sigma_1\) is fixed at the parent value during
amplification and changes at the time the splitting quantity moves into
the saturation regime.

For large \(N\),

\[
\varepsilon=\frac12-\frac{\sigma_2}{\sigma_1}
\]

decreases like \(1/N\). Figure 1 plots \(2(N-1)\varepsilon\) and shows
the temporal coincidence of the transition from the amplification regime
to the metastable regime with the spectral change.

\begin{figure}
\centering
\pandocbounded{\includegraphics[keepaspectratio,alt={Overlay of f and sigma2/sigma1 deviation}]{figs/cause_overlay_f_sigma_ratio_v1.png}}
\caption{Overlay of f and sigma2/sigma1 deviation}
\end{figure}

\emph{Figure 1: Splitting quantity \(f(\tau)\) and
\(2(N-1)(1/2-\sigma_2/\sigma_1)\) for \(N=5,40,300\). The dotted line
marks the time when the splitting quantity moves into the saturation
regime.}

Whether this indicator is the cause of the transition or a spectral
response accompanying it is undetermined. Here we treat it as an
observable that marks the transition time.

\begin{center}\rule{0.5\linewidth}{0.5pt}\end{center}

\subsection{\texorpdfstring{4. \(N\)-dependence of the metastable
amplitude}{4. N-dependence of the metastable amplitude}}\label{n-dependence-of-the-metastable-amplitude}

\subsubsection{4.1 Representative amplitude of relational
components}\label{representative-amplitude-of-relational-components}

The representative amplitude per relation in the metastable regime
follows

\[
A_{\mathrm{rel}}
\simeq
\frac{1}{\sqrt M}
=
\sqrt{\frac{2}{N(N-1)}}
\sim
\frac{\sqrt2}{N}.
\]

For \(N=3\)--\(300\), 18 values of \(N\), 38 runs in total, the mean of
\(A_{\mathrm{rel}}\sqrt M\) was \(1.00004\). A power fit for \(N\ge20\)
gave exponent \(\alpha=1.00820\), \(R^2=0.999990\), consistent with the
effective exponent \(1.00823\) of the exact binomial evaluated on the
same \(N\) grid. For \(N=300\), \(A_{\mathrm{rel}}=0.00472237\) against
\(1/\sqrt{44850}=0.00472192\).

\subsubsection{4.2 Convention and
dynamics}\label{convention-and-dynamics}

The \(1/\sqrt M\) law has two elements. First, from the convention
\(\sum_{e=1}^{M}|Z_e|^2=1\), perfect equipartition gives each component
amplitude \(1/\sqrt M\). Second, the actual dynamics approaches
\(\mathrm{PR}/M\to1\) in the metastable regime and narrows the amplitude
distribution in the edge basis. The nontrivial result is not the
normalization itself, but that the dynamics builds a
near-equipartitioned state in the edge basis.

Hence the observed \(1/N\) is not a fundamental law but the large-\(N\)
limit of \(1/\sqrt M\) with \(M=N(N-1)/2\).

\subsubsection{4.3 Relation to the spectral
ratio}\label{relation-to-the-spectral-ratio}

\[
\frac12-\frac{\sigma_2}{\sigma_1}=O\!\left(\frac1N\right)
\]

is also observed. That the \(1/\sqrt M\) law of the edge amplitude and
the \(1/N\) law of the spectral ratio arise from the same mechanism has
not yet been shown; their relation is a task for analysis of the
parent-generator spectrum.

\begin{center}\rule{0.5\linewidth}{0.5pt}\end{center}

\subsection{5. Inflow into orthogonal rotation
planes}\label{inflow-into-orthogonal-rotation-planes}

\subsubsection{5.1 Canonical decomposition of the antisymmetric
generator}\label{canonical-decomposition-of-the-antisymmetric-generator}

A real antisymmetric matrix \(K\) decomposes canonically into mutually
orthogonal two-dimensional rotation planes and a kernel {[}5{]}. In a
suitable orthogonal basis,

\[
K
\sim
\bigoplus_{j=1}^{r}
\begin{pmatrix}
0 & \sigma_j\\
-\sigma_j & 0
\end{pmatrix}
\oplus 0.
\]

Thus the basic unit of motion in this system is not a single axis but a
two-dimensional rotation plane.

For \(N=5\), the dominant-plane rotation rate was \(\sigma_1=3.742\) and
the next nonzero rate was \(\sigma_2=\sqrt2\). The cosine of the
principal angle between the two subspaces was zero within numerical
precision, confirming orthogonality.

\subsubsection{\texorpdfstring{5.2 The boundary between \(N=4\) and
\(N=5\)}{5.2 The boundary between N=4 and N=5}}\label{the-boundary-between-n4-and-n5}

We compare the number of nonzero rotation planes of the parent generator
with the splitting outcome.

{\def\LTcaptype{none} % do not increment counter
\begin{longtable}[]{@{}rrl@{}}
\toprule\noalign{}
N & second nonzero rate \(\sigma_2/\sigma_1\) & splitting outcome \\
\midrule\noalign{}
\endhead
\bottomrule\noalign{}
\endlastfoot
3 & 0.000 & no sustained expansion \\
4 & 0.000 & no sustained expansion \\
5 & 0.378 & expansion \\
6 & 0.386 & expansion \\
7 & 0.407 & expansion \\
\end{longtable}
}

For \(N\le4\) the parent generator has no second nonzero rotation
subspace; it first appears at \(N\ge5\). This boundary coincides with
the boundary of sustained splitting obtained in the earlier paper
{[}3{]}.

This result indicates that a nonzero rotation subspace outside the
parent plane may be necessary for sustained splitting. That
\(\sigma_2>0\) alone is a sufficient condition is not proven.

\subsubsection{5.3 Measurement of per-plane
norm}\label{measurement-of-per-plane-norm}

For the eigen-subspaces of the parent generator \(K(\arg v)\), we
projected the state norm at each time onto three groups:

\begin{enumerate}
\def\labelenumi{\arabic{enumi}.}
\tightlist
\item
  dominant plane \(P_1\): the 2D plane corresponding to \(\sigma_1\)
\item
  other rotation subspaces: \(0<\sigma<\sigma_1\)
\item
  kernel: \(\sigma=0\)
\end{enumerate}

From the definition of the splitting quantity via the parent-plane
basis, the identity

\[
\boxed{f(\tau)=1-E_{P_1}(\tau)}
\]

holds, where \(E_{P_1}\) is the projection norm ratio onto the dominant
plane. Thus \(f\) is not an abstract ``new wave quantity'' but the
fraction of conserved norm that has flowed out of the dominant plane
into its orthogonal complement.

\subsubsection{\texorpdfstring{5.4 Dense-matrix decomposition:
\(N=5,40\)}{5.4 Dense-matrix decomposition: N=5,40}}\label{dense-matrix-decomposition-n540}

For \(N=5\) and \(N=40\) we constructed the dense matrix \(K\) directly
and obtained the rotation subspaces and kernel from its
eigendecomposition. A clear spectral gap separated the nonzero
eigenvalue group from the numerical-zero group.

Projection dimensions:

{\def\LTcaptype{none} % do not increment counter
\begin{longtable}[]{@{}rrrrr@{}}
\toprule\noalign{}
N & M & dominant plane & other rotation subspaces & kernel \\
\midrule\noalign{}
\endhead
\bottomrule\noalign{}
\endlastfoot
5 & 10 & 2 & 4 & 4 \\
40 & 780 & 2 & 78 & 700 \\
\end{longtable}
}

The maximum deviation of the identity \(f=1-E_{P_1}\) was
\(3.7\times10^{-14}\) for \(N=5\) and \(3.8\times10^{-15}\) for
\(N=40\). The recorded \(f\) agreed bit-for-bit with the original run.

\begin{figure}
\centering
\pandocbounded{\includegraphics[keepaspectratio,alt={N=5 plane inflow}]{figs/planeflow_N00005_exact.png}}
\caption{N=5 plane inflow}
\end{figure}

\emph{Figure 2: \(N=5\). Projection norm ratios onto the dominant plane,
other rotation subspaces, and kernel. Outflow from the dominant plane
into the other two groups begins near \(\tau=1167\). Representative
metastable ratios are about 0.20, 0.33, 0.47.}

\begin{figure}
\centering
\pandocbounded{\includegraphics[keepaspectratio,alt={N=40 plane inflow}]{figs/planeflow_N00040_exact.png}}
\caption{N=40 plane inflow}
\end{figure}

\emph{Figure 3: \(N=40\). Outflow begins near \(\tau=2011\); in the
metastable regime the ratios are about 0.80 (dominant plane), 0.07
(other rotation subspaces), 0.13 (kernel).}

For both \(N=5\) and \(N=40\), the onset of the splitting quantity and
the norm outflow out of the dominant plane occur at the same time.

\subsubsection{\texorpdfstring{5.5 Low-rank decomposition:
\(N=300\)}{5.5 Low-rank decomposition: N=300}}\label{low-rank-decomposition-n300}

For \(N=300\), \(M=44850\), and dense-matrix eigendecomposition is not
feasible in this computing environment. We therefore used the low-rank
structure \(K=WJW^T\) and lifted the eigenpairs of \(JG\) to the edge
space to construct the projection subspaces.

To exclude numerical-zero eigenvalues, we set

\[
\sigma>10^{-6}\sigma_{\max}
\]

as the criterion for a nonzero rotation rate. Compared against the
dense-matrix method at \(N=40\), this approximation gave matching
subspace dimensions, with a maximum deviation of \(2.0\times10^{-15}\)
over all times and all projection groups.

For \(N=300\):

{\def\LTcaptype{none} % do not increment counter
\begin{longtable}[]{@{}lr@{}}
\toprule\noalign{}
group & dimension \\
\midrule\noalign{}
\endhead
\bottomrule\noalign{}
\endlastfoot
dominant plane & 2 \\
other rotation subspaces & 598 \\
kernel & 44250 \\
\end{longtable}
}

\begin{figure}
\centering
\pandocbounded{\includegraphics[keepaspectratio,alt={N=300 plane inflow}]{figs/planeflow_N00300_approx.png}}
\caption{N=300 plane inflow}
\end{figure}

\emph{Figure 4: \(N=300\). Low-rank method, nonzero threshold
\(10^{-6}\sigma_{\max}\). Outflow begins near \(\tau=4844\); in the
metastable regime the ratios are about 0.91 (dominant plane), 0.04
(other rotation subspaces), 0.05 (kernel).}

\subsubsection{5.6 Basis dependence}\label{basis-dependence}

In the edge basis the metastable state shows a near-equal-amplitude
distribution close to \(1/\sqrt M\). In the projection onto the
generator eigen-subspaces, the norm per dimension is not
equipartitioned. For example, at \(N=40\) the per-dimension
representative values are

\[
P_1:0.398,
\qquad
\text{other rotation}:0.0009,
\qquad
\text{kernel}:0.0002,
\]

so the dominant plane is strongly occupied. Therefore

\begin{itemize}
\tightlist
\item
  component equipartition in the edge basis
\item
  norm distribution in the generator eigen-subspaces
\end{itemize}

are different properties and cannot be identified.

\subsubsection{5.7 Conclusion}\label{conclusion-1}

From these observations, splitting can be described as the process

\[
\boxed{
\text{localization in a single dominant plane}
\rightarrow
\text{transverse amplification by phase feedback}
\rightarrow
\text{norm outflow into other rotation subspaces and the kernel}
\rightarrow
\text{bounded metastable distribution}
}.
\]

Whether the newly occupied subspaces correspond to real time, charge, or
internal degrees of freedom is not identified here.

\begin{center}\rule{0.5\linewidth}{0.5pt}\end{center}

\subsection{6. Physical identification of the unnamed directions, and
related
work}\label{physical-identification-of-the-unnamed-directions-and-related-work}

What this paper has shown is that orthogonal rotation subspaces and a
kernel distinct from the dominant plane exist, and that at splitting the
conserved norm flows into them. This existence, orthogonality, onset of
occupation, and norm inflow are derived and measured within the model;
they are not hypotheses.

Reading these unnamed directions as a time axis, internal charge axes,
or other physical quantities is a separate problem. Such identification
requires at least:

\begin{enumerate}
\def\labelenumi{\arabic{enumi}.}
\tightlist
\item
  a readout map that selects a fixed number of physical quantities from
  the \(N\)-dependent group of rotation planes,
\item
  a conservation or transformation law for each readout quantity,
\item
  an operational correspondence between the phase advance of the
  dominant plane and time measurement,
\item
  a correspondence between the internal rotational degrees of freedom
  and the sign/composition rules of charge.
\end{enumerate}

The task of a readout to a fixed number connects to the earlier paper
{[}4{]}, which showed that the number of waves is not an absolute
quantity within the system but depends on the readout resolution.

The isotropic two-planes obtained from the zero-square closure can be
compared with totally isotropic subspaces in pure-spinor and twistor
theory {[}6{]}. However, since the rotation planes here are constructed
from the time evolution of a state-dependent generator, this comparison
indicates a similarity of mathematical structure, not an identity.

As an example of emerging a time direction from rotational symmetry,
there is Walker's model obtaining an effective time axis from a gauged
rotational symmetry in four-dimensional Euclidean space {[}7{]}.
Constructions linking spacetime, internal symmetry, and Standard-Model
representations via \(D_4\), Clifford algebra, octonions, and \(E_8\)
have also been proposed {[}8--11{]}. These are comparison targets for
the physical identification we will examine in the future, but are not
used in this paper's proof of splitting arrest and new-direction
creation.

In a separate document we examine the hypothesis of reading

\[
(x,y,z,R)\oplus(Q_1,Q_2,Q_3,t)
\]

as two \(D_4\) structures. But this identification is not derived from
the numerical results of this paper alone. Here we treat the new
directions as unnamed and separate their physical naming to subsequent
research.

\begin{center}\rule{0.5\linewidth}{0.5pt}\end{center}

\subsection{7. Conclusion}\label{conclusion-2}

We measured the onset, saturation, and inflow destination of spontaneous
splitting in an \(N\)-body relational-wave closed system.

First, fixing the generator to the parent state produces no splitting;
exponential amplification occurs only when the generator is
reconstructed at each step from the state phases. The onset mechanism of
splitting is the phase feedback between the state and the generator.

Second, for large \(N\) the splitting quantity saturates while 80--92\%
of the norm remains in the parent plane. Hence the arrest is not a
simple depletion of parent norm. The mechanism determining the
\(N\)-dependent saturation value is not derived.

Third, the splitting quantity equals the fraction of norm that has
flowed out of the dominant plane, and the destination was the other
nonzero rotation subspaces and the generator kernel. Thus the
spontaneous splitting of this system can be formulated as a
redistribution of conserved norm among orthogonal subspaces.

Fourth, the metastable edge-component amplitude followed
\(A_{\mathrm{rel}}\simeq1/\sqrt M\), explained as the combination of the
total-norm convention and dynamical equipartition in the edge basis.

Therefore the single-wave state of this system is transversely unstable
under the state-dependent generator and moves conserved quantity into
new orthogonal subspaces. Metastabilization is not a cessation of wave
production but a bounded redistribution state that continues under a
finite outflow.

\begin{center}\rule{0.5\linewidth}{0.5pt}\end{center}

\subsection{8. Established results, hypotheses, and underived
items}\label{established-results-hypotheses-and-underived-items}

\subsubsection{8.1 Established}\label{established}

\begin{itemize}
\tightlist
\item
  The Cayley update preserves \(Z^\dagger Z\) and \(Z^T Z\).
\item
  Under a fixed generator the kernel seed does not amplify; under a
  state-dependent generator it does.
\item
  The initial splitting quantity grows exponentially.
\item
  The large-\(N\) saturation is not explained by depletion of
  parent-plane norm.
\item
  The metastable state is not a fixed point but a bounded state that
  keeps varying in time.
\item
  The edge-basis representative amplitude follows \(1/\sqrt M\).
\item
  For \(N\le4\) there is no second nonzero rotation plane; it appears at
  \(N\ge5\).
\item
  \(f=1-E_{P_1}\) holds; the splitting quantity is the outflow from the
  dominant plane.
\item
  For \(N=5,40\), dense-matrix eigendecomposition confirms the outflow
  destination is the other rotation subspaces and the kernel.
\item
  The \(N=300\) low-rank decomposition reproduces the dense-matrix
  method at \(N=40\) to machine precision.
\end{itemize}

\subsubsection{8.2 Observed but unproven
relations}\label{observed-but-unproven-relations}

\begin{itemize}
\tightlist
\item
  \(f_{\mathrm{sat}}(N)\) decreases with \(N\).
\item
  \(\sigma_2/\sigma_1\to1/2\) appears to hold.
\item
  \(1/2-\sigma_2/\sigma_1=O(1/N)\) appears to hold.
\item
  The existence of a second nonzero rotation plane may be a necessary
  condition for sustained splitting.
\end{itemize}

\subsubsection{8.3 Underived}\label{underived}

\begin{itemize}
\tightlist
\item
  Closed form of the saturation value \(f_{\mathrm{sat}}(N)\)
\item
  Condition under which the outflow stays bounded in the metastable
  regime
\item
  Analytic derivation of \(\sigma_2/\sigma_1\to1/2\)
\item
  Per-plane inflow within the other rotation subspaces
\item
  Dynamical role of the component flowing into the kernel
\item
  Physical readout of the unnamed rotation subspaces
\end{itemize}

\begin{center}\rule{0.5\linewidth}{0.5pt}\end{center}

\subsection{9. References}\label{references}

\subsubsection{Self-citations}\label{self-citations}

{[}1{]} N. Kihara, ``Linear upper bound of the generator rank and
three-direction saturation in N-body complete pairwise relational
waves,'' Zenodo, Concept DOI: 10.5281/zenodo.21465898, 2026.

{[}2{]} N. Kihara, ``Plane-decomposition readout in N-body
fixed-generator systems,'' Zenodo, Concept DOI: 10.5281/zenodo.21468959,
2026.

{[}3{]} N. Kihara, ``Onset and threefold classification of outcomes of
spontaneous splitting in N-body relational-wave closed systems,''
Zenodo, Concept DOI: 10.5281/zenodo.21486233, 2026.

{[}4{]} N. Kihara, ``The number of waves is the resolution of the
system,'' Zenodo, Concept DOI: 10.5281/zenodo.21486544, 2026.

\subsubsection{External references}\label{external-references}

{[}5{]} R. A. Horn and C. R. Johnson, \emph{Matrix Analysis}, 2nd ed.,
Cambridge University Press, 2013. Reference for the orthogonal canonical
form and 2D rotation-block decomposition of real antisymmetric matrices.

{[}6{]} A. Taghavi-Chabert, ``Twistor Geometry of Null Foliations in
Complex Euclidean Space,'' \emph{SIGMA} 13, 005 (2017). DOI:
10.3842/SIGMA.2017.005; arXiv:1505.06938. Structural comparison with
totally isotropic subspaces, pure spinors, and twistor space.

{[}7{]} M. L. Walker, ``Spontaneous Emergence of a Causal Time Axis in
Euclidean Space from a Gauged Rotational Symmetry Theory,''
\emph{Symmetry} 16(1), 4 (2024). DOI: 10.3390/sym16010004. Prior example
of emerging a time direction from rotational symmetry.

{[}8{]} F. D. T. Smith, ``Higgs and Fermions in D4-D5-E6 Model based on
Cl(0,8) Clifford Algebra,'' arXiv:hep-th/9403007 (1994). Prior proposal
connecting \(D_4\), Clifford algebra, and eight-dimensional structure to
physical representations.

{[}9{]} N. Furey, ``Three Generations, Two Unbroken Gauge Symmetries,
and One Eight-Dimensional Algebra,'' \emph{Physics Letters B} 785,
84--89 (2018); arXiv:1910.08395. Construction of three generations and
\(SU(3)\times U(1)\) from complex octonions.

{[}10{]} I. Todorov, ``Octonion Internal Space Algebra for the Standard
Model,'' \emph{Universe} 9(5), 222 (2023). DOI: 10.3390/universe9050222;
arXiv:2206.06912. Organization of the Standard-Model internal space via
octonions and Clifford algebra.

{[}11{]} C. A. Manogue, T. Dray, and R. A. Wilson, ``Octions: An \(E_8\)
Description of the Standard Model,'' \emph{Journal of Mathematical
Physics} 63, 081703 (2022). DOI: 10.1063/5.0095484; arXiv:2204.05310.
Comparison target for Standard-Model and Lorentz-algebra representations
within \(E_8\).

\begin{center}\rule{0.5\linewidth}{0.5pt}\end{center}

\subsection{10. Reproducibility}\label{reproducibility}

\subsubsection{10.1 Original programs}\label{original-programs}

\begin{itemize}
\tightlist
\item
  \texttt{run\_n\_scaling\_lowrank\_v1.py}
\item
  \texttt{run\_spontaneous\_splitting\_largeN\_v1.py}
\end{itemize}

The originals are fixed by SHA-256, and the instrumented programs do not
modify the time-evolution loop. In all instrumented runs, the splitting
quantity \(f\) agreed bit-for-bit with the original output.

\subsubsection{10.2 Instrumented programs}\label{instrumented-programs}

\begin{itemize}
\tightlist
\item
  \texttt{run\_metastable\_series\_v1.py}
\item
  \texttt{run\_cause\_instrumented\_v1.py}
\item
  \texttt{make\_cause\_overlay\_figure\_v1.py}
\item
  \texttt{run\_plane\_flow\_exact\_v1.py}
\item
  \texttt{run\_plane\_flow\_approx\_v1.py}
\item
  \texttt{make\_metastable\_series\_figure\_v1.py}
\item
  \texttt{trace\_parent\_iteration\_v1.py}
\end{itemize}

The nonzero-detection threshold of the low-rank approximation is
recorded in CLI arguments, output JSON, and the figures. Its default is
\(\texttt{sigma\_rel\_threshold}=10^{-6}\).

\subsubsection{10.3 Output data}\label{output-data}

\begin{itemize}
\tightlist
\item
  \texttt{cause\_instrumented\_result\_v1/}
\item
  \texttt{plane\_flow\_result\_v1/}
\item
  \texttt{metastable\_series\_result\_v1/}
\item
  \texttt{relational\_wave\_metastable\_amplitude\_N300\_result\_v1.md}
\item
  \texttt{parent\_crystal\_lowN\_observation\_record\_v1.md}
\end{itemize}

CSV files are regenerable outputs and are kept out of repository
tracking.

\begin{center}\rule{0.5\linewidth}{0.5pt}\end{center}

\subsection{11. Next experiments}\label{next-experiments}

\begin{enumerate}
\def\labelenumi{\arabic{enumi}.}
\item
  \textbf{Seed scan.} Scan \(\delta=10^{-9}\) to \(10^{-15}\) and
  measure the logarithmic shift of the saturation time and the
  independence of the saturation value.
\item
  \textbf{Transverse growth rate at large \(N\).} Measure the
  finite-time transverse growth rate simultaneously in runs including
  \(N=40, 300\).
\item
  \textbf{Inflow into individual rotation planes.} Decompose the ``other
  rotation subspaces'' into the individual \(\sigma_j\) planes and
  examine whether the inflow concentrates in a few planes or disperses
  over many.
\item
  \textbf{Saturation law.} Scan \(N\) densely to obtain the functional
  form of \(f_{\mathrm{sat}}(N)\).
\item
  \textbf{Role of the kernel component.} Measure the inflow into the
  kernel, its residence time, and the exchange rate with the rotation
  subspaces.
\end{enumerate}

The physical identification and the correspondence to \(D_4\) structures
and \(E_8\) diagonal gluing are separated from these dynamical
measurements and treated in subsequent papers.

\end{document}
